\documentclass[11pt]{article}
\usepackage[margin=1in]{geometry}
\usepackage{amsmath,amssymb,amsthm}
\usepackage{hyperref}
\usepackage{enumitem}
\usepackage{xurl}

\newtheorem{theorem}{Theorem}
\newtheorem{lemma}[theorem]{Lemma}
\newtheorem{corollary}[theorem]{Corollary}
\theoremstyle{definition}
\newtheorem{definition}[theorem]{Definition}
\theoremstyle{remark}
\newtheorem{remark}[theorem]{Remark}

\let\oldus\_
\renewcommand{\_}{\oldus\allowbreak}
\newcommand{\BadA}{\mathrm{Bad}_A}
\newcommand{\BadB}{\mathrm{Bad}_B}

\title{A disproof of Conjecture 00000000307\\[2pt]
\large $\mathrm{Bad}_A(i,j)=\emptyset$ for $i,j\le 0$ and $\mathrm{Bad}_B(i,j)=\emptyset$ for $i,j\le -3/4$}
\date{}

\begin{document}
\sloppy
\maketitle

\section{The conjecture}

Conjecture 00000000307 reads:
\begin{quote}
\emph{Definition:} Weighted bad approximation $\mathrm{Bad}(i,j)$ is the set of pairs simultaneously badly
approximable when the two coordinates are weighted by $q^{1+i}$ and $q^{1+j}$.
\emph{Conjecture:} For $i+j=0$, $\mathrm{Bad}(i,j)$ is a full-dimensional winning set; moreover, for every
weighting with $i+j<0$, the intersection remains full-dimensional, with winning given explicitly by a
Schmidt-game strategy.
\end{quote}
The Chinese version says the same (the second clause: for every choice of weights with $i+j<0$ the
intersection is still of full dimension).  There are two clauses:
\begin{enumerate}[label=(\Roman*)]
  \item for every $(i,j)$ with $i+j=0$, $\mathrm{Bad}(i,j)\subseteq\mathbb{R}^2$ has Hausdorff dimension $2$
        (and is winning);
  \item for every weighting $(i,j)$ with $i+j<0$, ``the intersection'' has Hausdorff dimension $2$ (and is
        winning).
\end{enumerate}
We only use the full-dimension conjuncts; the winning conjuncts are not formalized.

\section{Readings and conventions}

The phrase ``weighted by $q^{1+i}$'' admits two literal readings. In both, $q\ge 1$, $p$, $r$ range over
the integers, $i,j$ range over the reals (integer weights are a special case), $q^{1+i}$ is the real power,
and the plane carries its Euclidean metric.

\begin{definition}[Reading A: the weight multiplies the error $|x-p/q|$]
$(x,y)\in\BadA(i,j)$ iff there is $c>0$ such that for all integers $q\ge 1$, $p$, $r$:
\[ \max\bigl(q^{1+i}\,|x-p/q|,\ q^{1+j}\,|y-r/q|\bigr)\ \ge\ c. \]
\end{definition}

\begin{definition}[Reading B: the weight multiplies $\|qx\|$]
$(x,y)\in\BadB(i,j)$ iff there is $c>0$ such that for all integers $q\ge 1$:
$\max\bigl(q^{1+i}\|qx\|,\ q^{1+j}\|qy\|\bigr)\ge c$, where $\|t\|=\min_{p\in\mathbb{Z}}|t-p|$; equivalently,
for all integers $q\ge1$, $p$, $r$: $\max\bigl(q^{1+i}|qx-p|,\ q^{1+j}|qy-r|\bigr)\ge c$.
\end{definition}

The maximum is written in Lean as a disjunction of two inequalities. Variants that require both inequalities
(``each coordinate badly approximable''), a strict inequality, or all nonzero integers $q$ define subsets of
these sets, so they are empty as well.  ``The intersection'' in clause (II) is not defined in the text. We
only assume that, for the weighting $(i,j)$, it is a subset of $\mathrm{Bad}(i,j)$; this covers
$\bigcap_{i+j<0}\mathrm{Bad}(i,j)$ (the natural reading of the Chinese text) and $\mathrm{Bad}(i,j)\cap K$ for any $K$.
The Hausdorff dimension of the empty set is $0$ (Mathlib's \texttt{dimH\_empty}).

\section{The key lemma: two-dimensional Dirichlet}

\begin{lemma}[\texttt{dirichlet\_two}]\label{lem:dir}
For all real $x,y$ and every integer $N\ge1$ there are an integer $q$ with $1\le q\le N^2$ and integers $p,r$
with $|qx-p|<1/N$ and $|qy-r|<1/N$.
\end{lemma}
\begin{proof}
For $k=0,1,\dots,N^2$ put $f(k)=(\lfloor N\{kx\}\rfloor,\lfloor N\{ky\}\rfloor)\in\{0,\dots,N-1\}^2$, where
$\{t\}=t-\lfloor t\rfloor$. There are $N^2+1$ values of $k$ and $N^2$ boxes, so $f(k)=f(l)$ for some $l<k$
(pigeonhole). Then $|N\{kx\}-N\{lx\}|<1$, and with $q=k-l\in[1,N^2]$ and
$p=\lfloor kx\rfloor-\lfloor lx\rfloor$ we get $|qx-p|=|\{kx\}-\{lx\}|<1/N$; likewise for $y$
(Lean: \texttt{close\_of\_floor\_eq}).
\end{proof}

\section{The sets are empty}

\begin{theorem}[\texttt{badA\_eq\_empty}]
If $i\le 0$ and $j\le 0$ then $\BadA(i,j)=\emptyset$.
\end{theorem}
\begin{proof}
Let $c>0$ and pick $M\ge1$ with $1/M<c$. Lemma~\ref{lem:dir} with $N=M$ gives $q\ge1$, $p$, $r$ with
$|qx-p|,|qy-r|<1/M$. Since $q\ge1$ and $1+i\le 1$, $q^{1+i}\le q$, so
$q^{1+i}|x-p/q|\le q|x-p/q|=|qx-p|<1/M<c$, and likewise for $y$. So no $c$ works.
\end{proof}

\begin{theorem}[\texttt{badB\_eq\_empty}]
If $i\le -3/4$ and $j\le -3/4$ then $\BadB(i,j)=\emptyset$.
\end{theorem}
\begin{proof}
Let $c>0$ and pick $M\ge1$ with $1/M<c$. Lemma~\ref{lem:dir} with $N=M^2$ gives $1\le q\le M^4$, $p$, $r$ with
$|qx-p|,|qy-r|<1/M^2$. Since $1+i\le 1/4$ and $q\ge1$, $q^{1+i}\le q^{1/4}\le (M^4)^{1/4}=M$
(\texttt{rpow\_le\_of\_le\_pow\_four}), so $q^{1+i}|qx-p|<M\cdot M^{-2}=1/M<c$, and likewise for $y$.
\end{proof}

\begin{corollary}[the two clauses]
\begin{enumerate}[label=(\alph*)]
  \item Reading A, clause (I) fails: $\dim_H\BadA(0,0)=\dim_H\emptyset=0\neq2$, and $0+0=0$
        (\texttt{clause1\_false\_readingA}).
  \item Reading A, clause (II) fails: at the weighting $(-1/2,-1/2)$ (sum $-1<0$), every subset of
        $\BadA(-1/2,-1/2)=\emptyset$ has dimension $0\ne2$ (\texttt{clause2\_false\_readingA}); in particular
        $\dim_H\bigcap_{i+j<0}\BadA(i,j)=0$ (\texttt{dimH\_iInter\_badA}).
  \item Reading B, clause (II) fails: at $(-3/4,-3/4)$ (sum $-3/2<0$) every subset of
        $\BadB(-3/4,-3/4)=\emptyset$ has dimension $0$ (\texttt{clause2\_false\_readingB}); in particular
        $\dim_H\bigcap_{i+j<0}\BadB(i,j)=0$ (\texttt{dimH\_iInter\_badB}).
\end{enumerate}
Integer weightings are covered too: $\BadA(-1,-1)$ and $\BadB(-1,-1)$ are empty by the same theorems.
\end{corollary}

\begin{remark}[what is not refuted]
Under reading B, clause (I) is \emph{not} refuted: for example $\BadB(0,0)$ contains $\mathrm{Bad}\times\mathbb{R}$,
where $\mathrm{Bad}$ is the set of badly approximable reals. The conjecture is a conjunction, so it is false
under both readings, but under reading B only through clause (II). The winning conjuncts are not formalized
(an empty set is not winning, but we do not use this). We do not treat readings in which the weights enter as
exponents, or in which the defining inequality is only required for sufficiently large $q$.
\end{remark}

\section{Lean formalization}

File \texttt{lean/Conjecture307/Basic.lean} (238 lines, \texttt{import Mathlib} only, namespace
\texttt{Conjecture307}). The plane is \texttt{EuclideanSpace R (Fin 2)}; the sets are \texttt{BadA i j} and
\texttt{BadB i j} exactly as in Section 2. The main theorem is
\begin{verbatim}
theorem conjecture307_false :
    (not forall i j : R, i + j = 0 -> dimH (BadA i j) = 2) /\
    (forall I : R -> R -> Set Plane, (forall i j, I i j <= BadA i j) ->
      not forall i j : R, i + j < 0 -> dimH (I i j) = 2) /\
    (forall I : R -> R -> Set Plane, (forall i j, I i j <= BadB i j) ->
      not forall i j : R, i + j < 0 -> dimH (I i j) = 2)
\end{verbatim}
(written here in ASCII; in Lean $\subseteq$, $\forall$, $\neg$, $\mathbb{R}$). Supporting theorems:
\texttt{dirichlet\_two}, \texttt{close\_of\_floor\_eq}, \texttt{badA\_eq\_empty}, \texttt{badB\_eq\_empty},
\texttt{clause1\_false\_readingA}, \texttt{clause2\_false\_readingA}, \texttt{clause2\_false\_readingB},
\texttt{dimH\_iInter\_badA}, \texttt{dimH\_iInter\_badB}. The structure (empty set, then \texttt{dimH\_empty})
follows the accepted solution of conjecture 00000000310 in this repository (GPL-3.0).

\paragraph{Axiom audit.} \texttt{lake build} succeeds, and \texttt{\#print axioms} for
\texttt{conjecture307\_false}, \texttt{dimH\_iInter\_badA} and \texttt{dimH\_iInter\_badB} reports only
\texttt{[propext, Classical.choice, Quot.sound]}. There is no \texttt{sorry}, no \texttt{native\_decide} and no
added axiom.

\end{document}
